% ---
% title: Functional Analysis
% date: 2026/3/25
% tag: Math
% ---
These are my notes on functional analysis.

\section{Chap 0}

\subsection{Inner Product \& Outer Product}

\begin{center}
\begin{tabular}{llll}
\toprule
Operation & Scalar or Vector? & Measures & Physical Examples \\
\midrule
\textbf{Dot product} & Scalar & Alignment along a direction & Work, flux, projection \\
\textbf{Cross product} & Vector & Rotation / perpendicular effect & Torque, angular momentum, magnetic force \\
\bottomrule
\end{tabular}
\end{center}

\subsection{Distinct from Matrix Theory}

\begin{quote}
Functional analysis is the extension of finite-dimensional linear algebra (matrix theory) to infinite-dimensional spaces of functions and general operators.
\end{quote}

Once the dimension is infinite, several things change. The objects are functions rather than vectors in \(\mathbb{R}^n\): spaces like \(L^2\) or \(C([a,b])\). Linear maps become operators, often without a matrix representation and possibly unbounded. Facts from finite dimensions break: bounded no longer implies compact, norms are not equivalent, and bases behave differently. Convergence splits into norm, weak, and weak-* convergence. Eigenvalues become a spectrum, which can be continuous. Norms, completeness, and inner products turn into the main tools.

\subsection{Definition for InProduct in Function Space}

\(\mathbf{L}^2(-\pi,\pi)\) :

\[
(f,g)=\int_{-\pi}^{\pi}f(x)g(x)dx
\]

For Fourier series the orthogonal basis is the trigonometric functions.

\hrule

\subsection{Orthogonal Basis of Fourier Series}

For a function ( \(f(x)\) ) defined on ( \([-L, L]\) ):

\[
f(x) \sim \frac{a_0}{2} + \sum_{n=1}^{\infty} \big(a_n \cos\frac{n\pi x}{L} + b_n \sin\frac{n\pi x}{L}\big)
\]

The orthogonal basis functions are:

\[
{ 1, \cos\frac{n\pi x}{L}, \sin\frac{n\pi x}{L} }_{n=1}^{\infty}
\]

\hrule

\subsubsection{Orthogonality Relations}

Over the interval (\([-L, L]\)):

\[
\int_{-L}^{L} cos \frac{m\pi x}{L} cos \frac{n\pi x}{L},dx = 0 \quad (m \ne n)    \\ 
  \int_{-L}^{L} sin \frac{m\pi x}{L} sin \frac{n\pi x}{L},dx = 0 \quad (m \ne n)    \\
  \int_{-L}^{L} cos \frac{m\pi x}{L} sin \frac{n\pi x}{L},dx = 0 \quad (\forall m,n)
\]

The constant function (1) is orthogonal to all (\(\sin\)) and (\(\cos\)) functions for (\( n \ge 1\) ).

\hrule

\subsubsection{Intuition}

Each basis function is a “pure frequency”. Orthogonality means the projection of one basis onto another is zero: no overlap in “energy” between different frequencies. The Fourier coefficients \((a_n, b_n)\) are projections onto these directions, the same way a vector decomposes in Euclidean space.

\hrule

\subsubsection{Decomposition}

A function decomposes along the basis:

\[
f(x) \sim \sum_{k=0}^{n}(f,e_k)e_k
\]

\subsection{Fourier series convergence}

For a (2L)-periodic function (\(f(x)\)), the (N)-th partial sum is

\[
S_N(x) = \frac{a_0}{2} + \sum_{n=1}^{N} \left(a_n \cos \frac{n\pi x}{L} + b_n \sin \frac{n\pi x}{L}\right)
\]

It can also be written with the Dirichlet kernel (D\_N):

\[
S_N(x) = \int_{-L}^{L} f(y) D_N(x-y) , dy
\]

where

\[
D_N(t) = \frac{1}{2L} \frac{\sin\left((N+\frac12)\pi t / L\right)}{\sin(\pi t / 2L)}
\]

\hrule

\subsubsection{Convolution with Dirichlet kernel}

The Fourier partial sum is a convolution of (\(f\)) with (\(D_N\)):

\[
S_N(x) = \int_{-L}^{L} f(x - t) D_N(t) , dt
\]

(\(D_N(t)\)) is an approximate identity: as (\(N \to \infty\)), it behaves like a delta function, averaging symmetrically around (t=0).

\hrule

\subsubsection{Split around the point (x)}

Assume (f) is piecewise continuously differentiable, so the left and right limits exist:

\[
f(x^+) = \lim_{t \to 0^+} f(x+t), \quad f(x^-) = \lim_{t \to 0^-} f(x+t)
\]

Then we can write the integral as

\[
S_N(x) = \int_{-δ}^{δ} f(x-t) D_N(t) dt + \int_{|t|>δ} f(x-t) D_N(t) dt
\]

The second integral goes to 0 as (\(N \to \infty\)); the Dirichlet kernel decays away from \(0\). The first integral, around (t=0), gives the average of the left and right limits:

\[
\lim_{N \to \infty} \int_{-δ}^{δ} f(x-t) D_N(t) dt = \frac{f(x^+) + f(x^-)}{2}
\]

This uses the Dirichlet integral formula:

\[
\lim_{N \to \infty} \int_{-ε}^{ε} g(t) \frac{\sin((N+\frac12)t)}{\sin(t/2)} dt = \pi \cdot \frac{g(0^+) + g(0^-)}{2}
\]

\hrule

\subsubsection{收敛结果}

\[
\boxed{
\lim_{N \to \infty} S_N(x) = \frac{f(x^+) + f(x^-)}{2}
}
\]

If (f) is continuous at (x), this reduces to (f(x)). At a jump discontinuity, the Fourier series converges to the midpoint of the jump.

\hrule

Piecewise differentiability gives bounded variation, which controls the kernel integral. The Dirichlet kernel acts as an approximate delta, so the integral averages the two one-sided limits. This is the classical Dirichlet convergence theorem for Fourier series.

\subsection{Weak Convergence and Riemann integrals}

\subsubsection{Weak Convergence}

Let \((X, |\cdot|)\) be a Banach space, and let \({x_n} \subset X\). We say:

Weak convergence: \(x_n \rightharpoonup x\) if

\[
f(x_n) \to f(x) \quad \text{for all } f \in X^*,
\]

where \(X^*\) is the dual space (continuous linear functionals).

Weak-* convergence: \(x_n \stackrel{*}{\rightharpoonup} x\) in \(X^*\) if

\[
x_n(y) \to x(y) \quad \text{for all } y \in X
\]

This is common in spaces like \(L^\infty\) where the dual is bigger than the space itself.

\hrule

\subsubsection{Connection to Riemann Integrals}

Take \(f_n \in L^1([a,b])\). Weak convergence \(f_n \rightharpoonup f\) in \(L^1\) means

\[
\int_a^b f_n(x) g(x) dx \to \int_a^b f(x) g(x) dx \quad \forall g \in L^\infty([a,b])
\]

The \(f_n\) may not converge pointwise or in norm, but \(\int f_n g \to \int f g\) for every bounded test function. This generalizes the Riemann integral idea: the integral is the linear functional, and weak convergence is convergence under integration.

\hrule

\subsubsection{Intuition}

Strong convergence means functions converge in norm (the area under |f\_n - f| goes to 0). Weak convergence means they converge on average, through integration against test functions. Weak-* convergence is dual-level convergence, typical in spaces of measures or \(L^\infty\).

\hrule

\subsection{Weak convergence from finite dimensions}

In calculus, a sequence converges when:

\[
x_n \to x \iff |x_n - x| \to 0
\]

Extended to vectors (linear algebra), a sequence of vectors converges when every component converges, equivalently:

\[
\vec{v}_n \to \vec{v} \iff \|\vec{v}_n - \vec{v}\| \to 0
\]

This is strong convergence, the familiar kind.

\hrule

\subsubsection{What Weak Convergence Means}

Weak convergence does not require \(\|\vec{v}_n - \vec{v}\| \to 0\). It requires that measuring the sequence with any linear function gives converging values:

\begin{quote}
When you "measure" the sequence with any linear function, the measured values converge.
\end{quote}

In finite dimensions, every linear function is an inner product \(\langle a, x \rangle\), so weak convergence means:

\[
\vec{v}_n \rightharpoonup \vec{v} \iff \forall \vec{a},\quad \langle \vec{a}, \vec{v}_n \rangle \to \langle \vec{a}, \vec{v} \rangle
\]

In finite dimensions, weak and strong convergence coincide, so the distinction never comes up in linear algebra.

\hrule

\subsubsection{Why infinite dimensions are different}

Consider an infinite-dimensional setting, such as the space of infinite sequences. Take:

\[
e_1 = (1, 0, 0, 0, \ldots)
\]

\[
e_2 = (0, 1, 0, 0, \ldots)
\]

\[
e_3 = (0, 0, 1, 0, \ldots)
\]

Strongly, \(\|e_n - e_m\| = \sqrt{2}\) for all \(n \neq m\), so the vectors never get close.

Weakly, take any fixed \(a = (a_1, a_2, a_3, \ldots)\) and compute the inner product:

\[
\langle a, e_n \rangle = a_n \to 0
\]

because \(\sum a_i^2 < \infty\) forces \(a_n \to 0\). So:

\[
e_n \rightharpoonup 0 \quad \text{weakly, but } e_n \not\to 0 \text{ strongly}
\]

\hrule

\subsubsection{The two notions side by side}

\begin{center}
\begin{tabular}{ll}
\toprule
 & Intuition \\
\midrule
Strong convergence & The vectors themselves get closer and closer \\
Weak convergence & Every "viewpoint" sees them getting closer, but the vectors themselves can keep "running away" \\
\bottomrule
\end{tabular}
\end{center}

\hrule

\subsubsection{Where it shows up}

In probability theory, convergence in distribution, as in the Central Limit Theorem, is a form of weak convergence: it does not require the random variables to get close, only that the expected value of any continuous test function converges.

\subsection{Weak-* Convergence}

\hrule

\subsubsection{The dual space}

The dual space of a vector space \(X\) is

\[
X^* = \{\text{all linear functions } f : X \to \mathbb{R}\}
\]

From linear algebra: if \(X = \mathbb{R}^n\), every linear function on \(X\) looks like

\[
f(x) = a_1 x_1 + a_2 x_2 + \cdots + a_n x_n = \langle a, x \rangle
\]

So \(X^*\) is "another copy" of \(X\), made of row vectors acting on column vectors. In finite dimensions, \(X\) and \(X^*\) are indistinguishable. In infinite dimensions, they can be very different.

\hrule

\subsubsection{Levels of convergence}

Now take a sequence of linear functionals \(f_n \in X^*\) (not vectors).

\paragraph{Strong convergence in \(X^*\)}

\[
f_n \to f \iff \|f_n - f\|_{X^*} \to 0
\]

The functions themselves get close, in norm. Strictest requirement.

\paragraph{Weak convergence in \(X^*\)}

\[
f_n \rightharpoonup f \iff \forall g \in X^{**}, \quad g(f_n) \to g(f)
\]

Here \(X^{**}\) is the dual of the dual: you test \(f_n\) against all linear functions on \(X^*\). This is one level more abstract.

\paragraph{Weak-* convergence}

\[
f_n \overset{w^*}{\rightharpoonup} f \iff \forall x \in X, \quad f_n(x) \to f(x)
\]

You only test \(f_n\) against elements of the original space \(X\), not all of \(X^{**}\). This is the weakest requirement of the three.

\hrule

\subsubsection{A concrete example}

Let \(X = \ell^2\) (square-summable sequences), so \(X^* \cong \ell^2\) as well.

Take \(f_n = e_n^*\), the functional that picks out the \(n\)-th component:

\[
f_n(x) = x_n
\]

For weak-* convergence, take any fixed \(x = (x_1, x_2, \ldots) \in \ell^2\):

\[
f_n(x) = x_n \to 0
\]

since \(\sum x_i^2 < \infty\) forces \(x_n \to 0\). So \(f_n \overset{w^*}{\rightharpoonup} 0\).

Same conclusion as weak convergence here, because \(\ell^2\) is reflexive (explained below).

\hrule

\subsubsection{Reflexivity}

A space \(X\) is called reflexive if \(X^{**} = X\), i.e. the dual of the dual brings you back to where you started. In that case:

\[
\text{weak convergence} \iff \text{weak-* convergence}
\]

The gap appears in non-reflexive spaces:

\begin{center}
\begin{tabular}{lll}
\toprule
Space & Reflexive? & Weak \(\neq\) Weak-*? \\
\midrule
\(\mathbb{R}^n\) & Yes & No difference \\
\(\ell^2\), \(L^2\) & Yes & No difference \\
\(L^1\) & No & \textbf{Yes, they differ} \\
\(L^\infty\) & No & \textbf{Yes, they differ} \\
\(C([a,b])\) & No & \textbf{Yes, they differ} \\
\bottomrule
\end{tabular}
\end{center}

\hrule

\subsubsection{Implications}

\[
\text{Strong} \implies \text{Weak} \implies \text{Weak-*}
\]

Each arrow is strict in infinite-dimensional non-reflexive spaces. None of the arrows reverse in general.

\begin{center}
\begin{tabular}{lll}
\toprule
Convergence & What you test against & How strict \\
\midrule
Strong & Nothing — just the norm & Strictest \\
Weak & All of \(X^{**}\) & Middle \\
Weak-* & Only \(X\) itself & Weakest \\
\bottomrule
\end{tabular}
\end{center}

\hrule

\subsubsection{Banach-Alaoglu}

Banach-Alaoglu is what makes weak-* convergence useful:

\begin{quote}
The closed unit ball in \(X^*\) is always compact in the weak-* topology.
\end{quote}

In \(L^1\) or \(L^\infty\), strong compactness fails. But the closed unit ball is weak-* compact, so any bounded sequence of functionals has a weak-* convergent subsequence. This is the workhorse behind existence proofs in PDEs, optimization, and probability theory.

\section{无穷维空间的线性算子和谱分解}

\subsection{核心}

在 \(\mathbb{R}^n\) 中，线性算子就是矩阵，谱分解就是特征值分解。
进入无穷维空间后，"矩阵"变成了算子（如微分、积分、乘法算子），而"谱"的结构远比特征值丰富得多。

\begin{quote}
常用的空间：Hilbert 空间 \(H\)（完备内积空间），例如 \(L^2[0,1]\), \(L2[0,1]\), \(\ell^2(\mathbb{N})\), Sobolev 空间 \(H^1(\Omega)\)。
\end{quote}
